% ==============================================================
% Định dạng theo "Trình bày KLTN của chương trình đạt chuẩn quốc tế" — Trường ĐH
% Công nghệ, ĐHQGHN (docs/paper/Trình-bày-KLTN-của-chương-trình-đạt-chuẩn-quốc-tế.pdf).
% Mẫu không ghi thông số bằng chữ; các con số dưới đây đo trực tiếp từ file mẫu.
% Khóa luận dùng định dạng của mẫu nhưng viết bằng tiếng Việt.
% ==============================================================

% --- Trang & chữ ------------------------------------------------
% Lề: trái 1.5in, phải/trên/dưới 1in (đo từ mẫu). Khổ A4 theo quy định chung của trường
% (file mẫu là Letter; đổi a4paper→letterpaper ở đây và trong main.tex nếu cần đúng khổ mẫu).
\usepackage[a4paper, left=1.5in, right=1in, top=1in, bottom=1in,
            headheight=24pt, footskip=0.5in]{geometry}
\usepackage{fontspec}
\usepackage[english, main=vietnamese]{babel}
% Thân bài: Times New Roman (TeX Gyre Termes nếu máy không có).
\IfFontExistsTF{Times New Roman}
  {\setmainfont{Times New Roman}}
  {\setmainfont{texgyretermes}[Extension=.otf, UprightFont=*-regular,
     BoldFont=*-bold, ItalicFont=*-italic, BoldItalicFont=*-bolditalic]}
% Tiêu đề chương / danh mục / tài liệu tham khảo: Tahoma đậm (DejaVu Sans — cùng dáng — nếu không có).
\IfFontExistsTF{Tahoma}
  {\newfontfamily\headingfont{Tahoma}[Ligatures=TeX]}
  {\newfontfamily\headingfont{DejaVuSans}[Extension=.ttf, UprightFont=*, BoldFont=*-Bold, Ligatures=TeX]}
% 12pt, dãn dòng "1.5 lines" của Word (20.7pt); đoạn cách nhau 12pt, không thụt đầu dòng.
\usepackage{setspace}
\setstretch{1.43}
\setlength{\parindent}{0pt}
\setlength{\parskip}{12pt plus 2pt minus 1pt}
\setlength{\emergencystretch}{2em}  % giãn khoảng trắng thay vì tràn lề (đoạn văn, tài liệu tham khảo)

% --- Tiêu đề chương ----------------------------------------------
% Chương có số: "Chương 1" (Times 22pt, căn phải) rồi TÊN CHƯƠNG (Tahoma đậm 22pt, căn phải).
% Chương không số (Danh mục hình vẽ, DANH MỤC TỪ VIẾT TẮT, Tài liệu tham khảo): Tahoma đậm 20pt, căn giữa.
% Tên chương viết IN HOA ngay trong \chapter{...}, như mẫu.
\setkomafont{disposition}{\rmfamily\bfseries}
\setkomafont{chapter}{\normalfont}
\setkomafont{descriptionlabel}{\rmfamily\bfseries}
\renewcommand*{\chapterformat}{\chapapp~\thechapter}
\renewcommand*{\chapterlinesformat}[3]{%
  \ifblank{#2}
    {\centering\headingfont\bfseries\fontsize{20}{24}\selectfont #3\par}
    {\raggedleft{\rmfamily\mdseries\fontsize{22}{26}\selectfont #2\par}%
     \vspace{56pt}%
     {\headingfont\bfseries\fontsize{22}{26}\selectfont #3\par}}}
\RedeclareSectionCommand[beforeskip=40pt, afterskip=56pt]{chapter}

% --- Tiêu đề mục: 1.1. (16pt đậm), 1.1.1. (14pt đậm), 1.1.1.1. (13pt đậm nghiêng) ----
% Dấu chấm sau số mục (cả trong mục lục) do tùy chọn numbers=enddot của lớp tài liệu.
\setkomafont{section}{\fontsize{16}{19}\selectfont}
\setkomafont{subsection}{\fontsize{14}{17}\selectfont}
\setkomafont{subsubsection}{\fontsize{13}{16}\selectfont\itshape}
\RedeclareSectionCommands[beforeskip=18pt, afterskip=6pt]{section,subsection,subsubsection}
\setcounter{secnumdepth}{3}
\setcounter{tocdepth}{3}

% --- Mục lục ------------------------------------------------------
% Chương: chữ đậm, không ghi số, có dấu chấm dẫn; mục: "1.1."; tiểu mục cấp 3–4 in nghiêng.
\newcommand*{\tocgobble}[1]{}
\RedeclareSectionCommand[tocentrynumberformat=\tocgobble, tocnumwidth=0pt,
  tocentryformat=\rmfamily\bfseries, tocpagenumberformat=\rmfamily\bfseries,
  toclinefill=\TOCLineLeaderFill, tocbeforeskip=10pt]{chapter}
\RedeclareSectionCommand[tocnumwidth=2.8em]{section}
\RedeclareSectionCommand[tocnumwidth=3.6em]{subsection}
\RedeclareSectionCommand[tocnumwidth=4.6em]{subsubsection}
\RedeclareSectionCommands[tocentryformat=\itshape, tocpagenumberformat=\itshape]{subsection,subsubsection}

% --- Số trang: giữa, cuối trang ---------------------------------------
\usepackage[automark]{scrlayer-scrpage}
\clearpairofpagestyles
\cfoot*{\pagemark}
\setkomafont{pageheadfoot}{\normalfont}
\setkomafont{pagenumber}{\rmfamily\bfseries}

% --- Toán, hình, bảng, thuật toán ---------------------------------
\usepackage{amsmath, amssymb, amsthm, bm}
\usepackage{graphicx}
\graphicspath{{figures/}}
\usepackage{float, placeins}
\usepackage{pdflscape}  % bảng/hình khổ ngang: \begin{landscape}...\end{landscape}
% Chú thích: "Bảng 2.1 – ..." trên bảng, "Hình 2-1 – ..." dưới hình; đậm, 10pt, căn giữa.
% (Mẫu đánh số hình bằng gạch nối "2-1" và bảng bằng dấu chấm "2.1".)
\usepackage{caption}
\DeclareCaptionFont{tenpt}{\fontsize{10}{12}\selectfont}
\DeclareCaptionLabelSeparator{endash}{ -- }
\captionsetup{font={tenpt,bf}, labelsep=endash, justification=centering}
\captionsetup[table]{position=top}
% numbers=enddot không áp dụng cho nhãn hình/bảng: "Bảng 2.1 –", không phải "Bảng 2.1. –".
\renewcommand*{\figureformat}{\figurename~\thefigure}
\renewcommand*{\tableformat}{\tablename~\thetable}
\usepackage[font=tenpt, labelformat=parens, labelsep=space]{subcaption}
\usepackage{booktabs, multirow, makecell, tabularx, threeparttable, longtable}
\usepackage{siunitx}
\usepackage{algorithm}
\usepackage[noend]{algpseudocode}
\usepackage{xcolor}
\usepackage{tikz}
\usetikzlibrary{arrows.meta, positioning}
\usepackage{enumitem}
\setlist{nosep}
\setlist[itemize,enumerate]{leftmargin=*, labelindent=0.5in}

\numberwithin{equation}{chapter}
\numberwithin{figure}{chapter}
\numberwithin{table}{chapter}
\renewcommand{\thefigure}{\thechapter-\arabic{figure}}

% Danh mục hình vẽ / bảng biểu: "Hình 2-1 – Chú thích ........ 6"
\newcommand*{\loffigprefix}[1]{{\let\autodot\relax\figurename~#1~--\enskip}}
\newcommand*{\lottabprefix}[1]{{\let\autodot\relax\tablename~#1~--\enskip}}
\DeclareTOCStyleEntry[entrynumberformat=\loffigprefix, dynnumwidth=true, indent=0pt]{tocline}{figure}
\DeclareTOCStyleEntry[entrynumberformat=\lottabprefix, dynnumwidth=true, indent=0pt]{tocline}{table}

\addto\captionsvietnamese{%
  \renewcommand{\chaptername}{Chương}%
  \renewcommand{\appendixname}{Phụ lục}%
  \renewcommand{\figurename}{Hình}%
  \renewcommand{\tablename}{Bảng}%
  \renewcommand{\contentsname}{MỤC LỤC}%
  \renewcommand{\listfigurename}{Danh mục hình vẽ}%
  \renewcommand{\listtablename}{Danh mục bảng biểu}%
}
\floatname{algorithm}{Thuật toán}

% --- Trang phần đầu (LỜI CAM ĐOAN, TÓM TẮT, MỤC LỤC…): Times đậm 16pt, căn giữa
\newcommand{\frontheading}[1]{%
  {\centering\rmfamily\bfseries\fontsize{16}{20}\selectfont #1\par}\vspace{12pt}}
\newcommand{\frontpage}[1]{\clearpage\frontheading{#1}}
\deftocheading{toc}{\clearpage\frontheading{#1}}


% --- Tài liệu tham khảo: IEEE, đánh số theo thứ tự trích dẫn -----
\usepackage[autostyle=false, style=english]{csquotes}  % csquotes chưa có kiểu tiếng Việt
\usepackage[backend=biber, style=ieee, sorting=none]{biblatex}
\addbibresource{references.bib}
\DeclareLanguageMapping{vietnamese}{english}  % biblatex chưa có bản tiếng Việt: các chữ "in", "vol.", "pp." giữ kiểu IEEE
\renewrobustcmd*{\bibinitdelim}{}  % "F.T. Ulaby" như mẫu

% --- Liên kết (để cuối) -----------------------------------------
\usepackage[hidelinks, unicode]{hyperref}

% --- Lệnh tiện dụng ---------------------------------------------
\newcommand{\todo}[1]{\textcolor{red}{[TODO: #1]}}
